% ---------------------------------------------------------------------------
% Section 3: Theory. Draft v1, 2026-09-11 (Day 2).
%
% Every displayed equation is annotated with the file that generates the
% corresponding number. DAY 11 VERIFICATION: re-check each annotated line
% against vendor/gleam/{theorem2_learning,lstd_lambda,bias_stability,
% kernel_general}.py and prereg_e7_bias.py before submission. The formulas
% here were transcribed from the frozen prereg_e7_bias.py docstring (S index)
% and from kernel_general.py (critic, gradient), which reproduces the frozen
% prereg_e10.py tables to relative 6e-7 at all eight grid points.
%
% This section absorbs what were sections 2.3 and 2.4 in the earlier draft.
% ---------------------------------------------------------------------------

\section{Theory: the aliased equilibrium map and two thresholds}
\label{sec:theory}

Everything the experiments test is a closed-form consequence of the model of
\S\ref{sec:testbed} under one restriction: the critic sees only $v$. For a
linear policy $u=-kv$ the closed loop $F_k$ has stationary covariance
$\Sigma(k)$ from a discrete Lyapunov equation and deployed cost
$J(k)=(q+rk^2)\Sigma_{11}(k)$; from it, the best memoryless gain
$k_{\mathrm{ml}}=\arg\min_k J$, the full-state optimum $J^\ast$ (Riccati),
and the large-gain plateau $J\to r\Theta K(0)$. Derivations are in
Appendix~\ref{app:theory}; the generators are released.
% generator: theorem2_learning.J_deploy; gleam_bench.J_gain

\subsection{The aliased critic}

A critic restricted to the observation is $\hat V(v) = \theta_0 + \theta_1 v +
\theta_2 v^2$ with $\phi = (1, v, v^2)$. Its LSTD($\lambda$) fixed point solves
$A\theta = b$ with
\begin{equation}
A = \sum_{j\ge 0} (\beta\lambda)^j\,\mathbb{E}\!\left[\phi_{t-j}\,(\phi_t - \beta\phi_{t+1})^\top\right],
\qquad
b = \sum_{j\ge 0} (\beta\lambda)^j\,\mathbb{E}\!\left[\phi_{t-j}\, c_t\right].
\label{eq:lstd}
\end{equation}
% generator: lstd_lambda.lstd_lambda_theta2; kernel_general.theta2
By parity $\theta_1 = 0$, and with $x_j := e_1^\top F_k^{\,j}\,\Sigma\, e_1$
the lag-$j$ autocovariance of $v$, the curvature coefficient is
\begin{equation}
\theta_2(\lambda,k) \;=\; \frac{h\,(q + r k^2)\,\sum_{j\ge 0} (\beta\lambda)^j x_j^2}
{\sum_{j\ge 0} (\beta\lambda)^j\,\bigl(x_j^2 - \beta\, x_{j+1}^2\bigr)}.
\label{eq:theta2}
\end{equation}
This is the whole source of the confound. The critic is fed the persistence of
$v^2$ across lags, which under aliasing is the hidden bath's persistence
wearing the observation's clothes; the bootstrap weight $(\beta\lambda)^j$
decides how many lags it is allowed to see.

\subsection{The equilibrium map}

For the exploring policy $u = -kv + \eta$, $\eta\sim\mathcal N(0,\sigma_e^2)$,
the expected GAE($\lambda$) policy gradient with score
$s_t = -v_t\eta_t/\sigma_e^2$ is
\begin{equation}
\hat g(k;\lambda) \;=\; \mathbb{E}\Bigl[\, s_t \sum_{j\ge 0} (\beta\lambda)^j\,\delta_{t+j} \Bigr],
\qquad
\delta_t = c_t + \beta \hat V(v_{t+1}) - \hat V(v_t).
\label{eq:ghat}
\end{equation}
% generator: lstd_lambda.gae_gradient; kernel_general.ghat
Because the closed loop is Gaussian, every term is an Isserlis sum of lag
moments, and the infinite bootstrap series collapse to matrix resolvents ---
$\sum_j (\beta\lambda)^j x_j = e_1^\top (I - \beta\lambda F_k)^{-1}\Sigma e_1$,
and its quadratic analogue through $F_k\!\otimes\!F_k$ --- with no truncation
anywhere (Appendix~\ref{app:theory}). The \emph{aliased equilibrium map}
\begin{equation}
k_{\mathrm{AC}}(\lambda) \;:=\; \text{the largest root of } \hat g(\,\cdot\,;\lambda)
\label{eq:kAC}
\end{equation}
is the fixed point of the averaged actor--critic update. We say fixed point
rather than limit: whether training arrives there is one of the things the
experiments measure. At $\lambda = 0$ every $j\ge1$ term vanishes and
\eqref{eq:theta2}--\eqref{eq:kAC} reduce to the TD(0) closed form of the
companion derivation. On the reference instance $k_{\mathrm{AC}}$ runs from
$24.79$ at $\lambda=0$ to $1.507$ at $\lambda=1$, against
$k_{\mathrm{ml}} = 1.649$: the learning target of the memoryless baseline moves
by an order of magnitude under a hyperparameter that is rarely reported to two
decimals, and it moves \emph{away} from the true memoryless optimum as the
horizon shortens.
% numbers: prereg_e10.py C["GLEAM_g1_h0.05"]; gleam_bench.ac_equilibrium (lambda=0)

\subsection{The cost law and its threshold}

The learner's excess cost is $E(\lambda) := J(k_{\mathrm{AC}}(\lambda))/J_{\mathrm{ml}}$.
It sits on a plateau near $r\Theta K(0)/J_{\mathrm{ml}}$ while the bootstrap
horizon $(1-\beta\lambda)^{-1}$ is short against the bath's correlation time,
and collapses to $1$ once the horizon exceeds it. The \emph{cost-law threshold}
$\lambda_{1/2}$ is where $E$ has fallen halfway from plateau to floor. On the
four instances it is $0.9966, 0.9854, 0.9580, 0.8731$ for
$\gamma h = 0.02, 0.05, 0.10, 0.20$ --- monotone in the ratio of horizon to
correlation time, as it must be. Because experiments sample $\lambda$ on an
eight-point grid we also freeze the \emph{grid-matched} value --- the same
midpoint estimator applied to the exact $E$ on that grid --- so measurement
error and interpolation bias are never conflated; and the estimator is
invariant under scaling and shifting the profile, so it needs no reference
value (\S\ref{sec:results}).
% numbers: prereg_phase2.LAMBDA_HALF; prereg_e10.py GRID_LAMBDA_HALF

\subsection{Off-manifold stability and the second threshold}

Whether training can hold $k_{\mathrm{AC}}$ is a separate question from where
it is. Give the policy one direction off the gain manifold, $u = -kv + b$. By
parity $b=0$ is a fixed point of the averaged update for $b$; its stability is
exact. With $\mu_v = e_1^\top(I - F_k)^{-1}G$,
$\alpha = q\mu_v - kr(1 - k\mu_v)$,
$\Phi = \beta\lambda\, e_1^\top (I - \beta\lambda F_k)^{-1} G$,
$X = \sum_l (\beta\lambda)^l x_l$, $X_1 = \sum_l (\beta\lambda)^l x_{l+1}$, and
$\kappa = 2h\alpha X/(X - \beta X_1)$,
\begin{equation}
S(\lambda,k) \;=\; 2hr\,(1 - k\mu_v) \;+\; 2h\alpha\,\Phi \;+\; \kappa\,\Phi\,\frac{1-\lambda}{\lambda},
\qquad
\frac{\partial \hat g_b}{\partial b}\Big|_{b=0} = -S,
\label{eq:S}
\end{equation}
% generator: bias_stability.bias_stability; frozen in prereg_e7_bias.py
so $b=0$ is stable iff $S>0$. The first two terms are the direct restoring
force. The third is the aliased critic's contribution, and it carries the
factor $(1-\lambda)/\lambda$ --- the bootstrap weight itself. It vanishes at
$\lambda=1$ and grows without bound as $\lambda$ falls; below a critical value
it overwhelms the restoring terms and the policy leaves the manifold. The
\emph{stability threshold} $\lambda^\ast$ is the sign change of
$S(\lambda, k_{\mathrm{AC}}(\lambda))$: $0.9956, 0.9759, 0.9010$ on the first
three instances and no crossing on the fourth, where the equilibrium is stable
at every $\lambda$. Below $\lambda^\ast$ the theory predicts that training does
not converge and its measured quality depends on when it is stopped. The closed form agrees with a Monte Carlo central difference of
$\mathbb{E}[\hat g_b]$ (critic refit at each $b$) within 1--3\% at ten points
spanning both signs of $S$ and three instances.
% numbers: prereg_e7_bias.PREDICTIONS; MC table in prereg_e7_bias.py docstring

\subsection{Scope}

\begin{proposition}[Scope]
If the transition law of the hidden state does not depend on the action, then
for every $\lambda$ and every observation-restricted critic $\hat V$, the
critic's bias cancels identically from the expected policy gradient.
\end{proposition}
\emph{Proof.} The advantage is built from
$\delta_t = c_t + \beta\hat V(o_{t+1}) - \hat V(o_t)$, and the action enters
the expected advantage of $a$ at $o_t$ only through
$\mathbb{E}[\hat V(o_{t+1})\mid o_t, a]$. If transitions are
action-independent this term is the same for every $a$, so it is a baseline
and drops from every action preference.~$\square$

The proposition says where the confound is \emph{absent}. Displacing an
equilibrium needs more than non-cancellation: it needs a hidden coordinate the
action can \emph{steer}. Environments where actions merely \emph{reveal} a
frozen latent satisfy neither hypothesis, and the theory makes no prediction
for them. \S\ref{sec:results} tests both sides.

\subsection{What is a property of the family and what is not}

None of the formulas above uses the state dimension: they hold for any Prony
kernel $K(\tau)=\sum_i c_i e^{-\gamma_i\tau}$, and \S\ref{sec:results}
recomputes them at dimensions 2--4. Two things do \emph{not} carry across
kernel shapes --- a single timescale ordering $\lambda_{1/2}$, and the ordering
$\lambda^\ast<\lambda_{1/2}$ --- and we report both. The companion derivation
adds that, at matched autocovariance, the equilibrium \emph{location} moves
$3.4\times$ across Gaussian, bounded and kurtotic drives while the \emph{cost}
at equilibrium stays within 1--7\%: thresholds and cost levels are robust,
locations are fragile. That split predicts this paper's own ledger.
% numbers: kernel-generality-findings.md; GLEAM lemma-a-findings.md
